% !TeX root = ../main.tex
\section{Conclusion}\label{sec:conclusion}

This paper traced a single arc: from the arithmetic of cards to agents that play. The probabilistic core (Section~\ref{sec:probability}) is elementary but not optional---pot odds, the $1/\sqrt{n}$ wall of Monte Carlo estimation, and the variance computation behind \eqref{eq:handsneeded} set the boundaries within which every poker decision system must operate. The game-theoretic core (Sections~\ref{sec:gametheory}--\ref{sec:cfr}) supplies the conceptual shift that separates poker AI from classical game AI: in imperfect-information games, safety means equilibrium, equilibrium means mixed strategies, and mixed strategies mean accepting negative-EV actions (bluffs, calls that lose on average) as the premium paid for being unreadable. Kuhn poker compresses the entire lesson into twelve information sets, and its exact bluff--value arithmetic---one bluff per three value bets, matching the caller's pot odds to four decimal places---is the same arithmetic that governs no-limit bet sizing at any depth of stack.

The algorithmic ladder then formalised a trade-off that the experimental study made concrete. Rule-based expert systems (Section~\ref{sec:evaluation}) are transparent and fast but statically exploitable; Monte Carlo agents sharpen estimation while inheriting static decision thresholds; equilibrium blueprints computed by sampled CFR (Section~\ref{sec:cfr}) trade exploitability for unexploitability at the cost of an abstraction's fidelity; function approximation (Deep CFR) and trial-and-error learning (Q-learning, Section~\ref{sec:rl}) relax the tabular bottleneck in exchange for approximation error and non-stationarity hazards. The 100{,}000-hand tournament of Section~\ref{sec:experiments} measured where these trade-offs land when five paradigms share one table, under conditions---fixed stacks, rotating positions, seeded deals---designed so that observed differences are attributable to the agents rather than to seat luck, wealth compounding, or measurement artifacts.

Three engineering lessons generalise beyond poker. First, \emph{measurement discipline precedes interpretation}: our own investigation began with an apparent positional bias that dissolved under correct accounting, a reminder that in high-variance domains a plausible statistic unaccompanied by a variance calculation is indistinguishable from noise. Second, \emph{abstraction is not merely compression but a design language}: the choice of buckets, bet sizes, and betting-pattern keys encodes hypotheses about which strategic distinctions matter, and a shared abstraction is what makes controlled comparisons between paradigms meaningful at all. Third, \emph{guarantees and performance are different currencies}: the equilibrium-flavoured agents are certifiably hard to exploit within their abstract game, yet the tournament's winner is whichever agent's inductive biases happen to fit the field---a separation that mirrors the exploitation-versus-safety tension the theory predicts, and that future systems reconcile by combining both, as the superhuman programs do with blueprint equilibrium plus targeted real-time search.

The natural extensions of this work follow the field's trajectory: deeper stacks and finer action abstractions; exploitability measurement of the trained blueprints against best responses computed in the abstract game; end-to-end deep RL with NFSP-style averaging to place the learning paradigm on an equal representational footing with Deep CFR; and opponent modelling layered on top of equilibrium baselines, which is where commercial poker tools now concentrate their effort. The game has been kind to its students: each time the community declared the frontier closed---limit solved, heads-up solved, multiplayer superhuman---the remaining imperfections proved interesting enough to open the next one. Card probabilities, regret, and reinforcement learning will keep finding work at that frontier.
